% Options for packages loaded elsewhere
\PassOptionsToPackage{unicode}{hyperref}
\PassOptionsToPackage{hyphens}{url}
\documentclass[
  11pt,
]{article}
\usepackage{xcolor}
\usepackage[margin=2.5cm]{geometry}
\usepackage{amsmath,amssymb}
\setcounter{secnumdepth}{5}
\usepackage{iftex}
\ifPDFTeX
  \usepackage[T1]{fontenc}
  \usepackage[utf8]{inputenc}
  \usepackage{textcomp} % provide euro and other symbols
\else % if luatex or xetex
  \usepackage{unicode-math} % this also loads fontspec
  \defaultfontfeatures{Scale=MatchLowercase}
  \defaultfontfeatures[\rmfamily]{Ligatures=TeX,Scale=1}
\fi
\usepackage{lmodern}
\ifPDFTeX\else
  % xetex/luatex font selection
\fi
% Use upquote if available, for straight quotes in verbatim environments
\IfFileExists{upquote.sty}{\usepackage{upquote}}{}
\IfFileExists{microtype.sty}{% use microtype if available
  \usepackage[]{microtype}
  \UseMicrotypeSet[protrusion]{basicmath} % disable protrusion for tt fonts
}{}
\makeatletter
\@ifundefined{KOMAClassName}{% if non-KOMA class
  \IfFileExists{parskip.sty}{%
    \usepackage{parskip}
  }{% else
    \setlength{\parindent}{0pt}
    \setlength{\parskip}{6pt plus 2pt minus 1pt}}
}{% if KOMA class
  \KOMAoptions{parskip=half}}
\makeatother
\usepackage{color}
\usepackage{fancyvrb}
\newcommand{\VerbBar}{|}
\newcommand{\VERB}{\Verb[commandchars=\\\{\}]}
\DefineVerbatimEnvironment{Highlighting}{Verbatim}{commandchars=\\\{\}}
% Add ',fontsize=\small' for more characters per line
\usepackage{framed}
\definecolor{shadecolor}{RGB}{248,248,248}
\newenvironment{Shaded}{\begin{snugshade}}{\end{snugshade}}
\newcommand{\AlertTok}[1]{\textcolor[rgb]{0.94,0.16,0.16}{#1}}
\newcommand{\AnnotationTok}[1]{\textcolor[rgb]{0.56,0.35,0.01}{\textbf{\textit{#1}}}}
\newcommand{\AttributeTok}[1]{\textcolor[rgb]{0.13,0.29,0.53}{#1}}
\newcommand{\BaseNTok}[1]{\textcolor[rgb]{0.00,0.00,0.81}{#1}}
\newcommand{\BuiltInTok}[1]{#1}
\newcommand{\CharTok}[1]{\textcolor[rgb]{0.31,0.60,0.02}{#1}}
\newcommand{\CommentTok}[1]{\textcolor[rgb]{0.56,0.35,0.01}{\textit{#1}}}
\newcommand{\CommentVarTok}[1]{\textcolor[rgb]{0.56,0.35,0.01}{\textbf{\textit{#1}}}}
\newcommand{\ConstantTok}[1]{\textcolor[rgb]{0.56,0.35,0.01}{#1}}
\newcommand{\ControlFlowTok}[1]{\textcolor[rgb]{0.13,0.29,0.53}{\textbf{#1}}}
\newcommand{\DataTypeTok}[1]{\textcolor[rgb]{0.13,0.29,0.53}{#1}}
\newcommand{\DecValTok}[1]{\textcolor[rgb]{0.00,0.00,0.81}{#1}}
\newcommand{\DocumentationTok}[1]{\textcolor[rgb]{0.56,0.35,0.01}{\textbf{\textit{#1}}}}
\newcommand{\ErrorTok}[1]{\textcolor[rgb]{0.64,0.00,0.00}{\textbf{#1}}}
\newcommand{\ExtensionTok}[1]{#1}
\newcommand{\FloatTok}[1]{\textcolor[rgb]{0.00,0.00,0.81}{#1}}
\newcommand{\FunctionTok}[1]{\textcolor[rgb]{0.13,0.29,0.53}{\textbf{#1}}}
\newcommand{\ImportTok}[1]{#1}
\newcommand{\InformationTok}[1]{\textcolor[rgb]{0.56,0.35,0.01}{\textbf{\textit{#1}}}}
\newcommand{\KeywordTok}[1]{\textcolor[rgb]{0.13,0.29,0.53}{\textbf{#1}}}
\newcommand{\NormalTok}[1]{#1}
\newcommand{\OperatorTok}[1]{\textcolor[rgb]{0.81,0.36,0.00}{\textbf{#1}}}
\newcommand{\OtherTok}[1]{\textcolor[rgb]{0.56,0.35,0.01}{#1}}
\newcommand{\PreprocessorTok}[1]{\textcolor[rgb]{0.56,0.35,0.01}{\textit{#1}}}
\newcommand{\RegionMarkerTok}[1]{#1}
\newcommand{\SpecialCharTok}[1]{\textcolor[rgb]{0.81,0.36,0.00}{\textbf{#1}}}
\newcommand{\SpecialStringTok}[1]{\textcolor[rgb]{0.31,0.60,0.02}{#1}}
\newcommand{\StringTok}[1]{\textcolor[rgb]{0.31,0.60,0.02}{#1}}
\newcommand{\VariableTok}[1]{\textcolor[rgb]{0.00,0.00,0.00}{#1}}
\newcommand{\VerbatimStringTok}[1]{\textcolor[rgb]{0.31,0.60,0.02}{#1}}
\newcommand{\WarningTok}[1]{\textcolor[rgb]{0.56,0.35,0.01}{\textbf{\textit{#1}}}}
\usepackage{longtable,booktabs,array}
\newcounter{none} % for unnumbered tables
\usepackage{calc} % for calculating minipage widths
% Correct order of tables after \paragraph or \subparagraph
\usepackage{etoolbox}
\makeatletter
\patchcmd\longtable{\par}{\if@noskipsec\mbox{}\fi\par}{}{}
\makeatother
% Allow footnotes in longtable head/foot
\IfFileExists{footnotehyper.sty}{\usepackage{footnotehyper}}{\usepackage{footnote}}
\makesavenoteenv{longtable}
\usepackage{graphicx}
\makeatletter
\newsavebox\pandoc@box
\newcommand*\pandocbounded[1]{% scales image to fit in text height/width
  \sbox\pandoc@box{#1}%
  \Gscale@div\@tempa{\textheight}{\dimexpr\ht\pandoc@box+\dp\pandoc@box\relax}%
  \Gscale@div\@tempb{\linewidth}{\wd\pandoc@box}%
  \ifdim\@tempb\p@<\@tempa\p@\let\@tempa\@tempb\fi% select the smaller of both
  \ifdim\@tempa\p@<\p@\scalebox{\@tempa}{\usebox\pandoc@box}%
  \else\usebox{\pandoc@box}%
  \fi%
}
% Set default figure placement to htbp
\def\fps@figure{htbp}
\makeatother
% definitions for citeproc citations
\NewDocumentCommand\citeproctext{}{}
\NewDocumentCommand\citeproc{mm}{%
  \begingroup\def\citeproctext{#2}\cite{#1}\endgroup}
\makeatletter
 % allow citations to break across lines
 \let\@cite@ofmt\@firstofone
 % avoid brackets around text for \cite:
 \def\@biblabel#1{}
 \def\@cite#1#2{{#1\if@tempswa , #2\fi}}
\makeatother
\newlength{\cslhangindent}
\setlength{\cslhangindent}{1.5em}
\newlength{\csllabelwidth}
\setlength{\csllabelwidth}{3em}
\newenvironment{CSLReferences}[2] % #1 hanging-indent, #2 entry-spacing
 {\begin{list}{}{%
  \setlength{\itemindent}{0pt}
  \setlength{\leftmargin}{0pt}
  \setlength{\parsep}{0pt}
  % turn on hanging indent if param 1 is 1
  \ifodd #1
   \setlength{\leftmargin}{\cslhangindent}
   \setlength{\itemindent}{-1\cslhangindent}
  \fi
  % set entry spacing
  \setlength{\itemsep}{#2\baselineskip}}}
 {\end{list}}
\usepackage{calc}
\newcommand{\CSLBlock}[1]{\hfill\break\parbox[t]{\linewidth}{\strut\ignorespaces#1\strut}}
\newcommand{\CSLLeftMargin}[1]{\parbox[t]{\csllabelwidth}{\strut#1\strut}}
\newcommand{\CSLRightInline}[1]{\parbox[t]{\linewidth - \csllabelwidth}{\strut#1\strut}}
\newcommand{\CSLIndent}[1]{\hspace{\cslhangindent}#1}
\setlength{\emergencystretch}{3em} % prevent overfull lines
\providecommand{\tightlist}{%
  \setlength{\itemsep}{0pt}\setlength{\parskip}{0pt}}
\usepackage{booktabs}
\usepackage{longtable}
\usepackage{float}
\usepackage{amsmath}
\usepackage{amssymb}
\floatplacement{figure}{H}
\floatplacement{table}{H}
\usepackage{bookmark}
\IfFileExists{xurl.sty}{\usepackage{xurl}}{} % add URL line breaks if available
\urlstyle{same}
\hypersetup{
  pdftitle={gretaR: A Torch-Based Probabilistic Programming Framework for Bayesian Statistical Modelling in R},
  hidelinks,
  pdfcreator={LaTeX via pandoc}}

\title{gretaR: A Torch-Based Probabilistic Programming Framework for
Bayesian Statistical Modelling in R}
\usepackage{etoolbox}
\makeatletter
\providecommand{\subtitle}[1]{% add subtitle to \maketitle
  \apptocmd{\@title}{\par {\large #1 \par}}{}{}
}
\makeatother
\subtitle{Technical Documentation and Reference}
\author{true}
\date{11 April 2026}

\begin{document}
\maketitle
\begin{abstract}
gretaR is an R package for Bayesian statistical modelling that provides
an R-native domain-specific language (DSL) for defining probabilistic
models, compiled to differentiable computation graphs for inference via
Hamiltonian Monte Carlo, the No-U-Turn Sampler, variational inference,
and maximum a posteriori estimation. The package offers a dual-backend
architecture: a native torch backend for GPU-capable custom models and
an automatically generated Stan backend for production-speed inference
on standard models. This document provides comprehensive technical
documentation of the package architecture, algorithms, and
implementation.
\end{abstract}

{
\setcounter{tocdepth}{3}
\tableofcontents
}
\section{Introduction}\label{introduction}

gretaR is a probabilistic programming package for R that enables
Bayesian statistical modelling using native R syntax. Models are defined
interactively by manipulating \texttt{gretaR\_array} objects --- data
wrappers, random variables, and deterministic transformations --- which
lazily construct a directed acyclic graph (DAG). This graph is then
compiled into a differentiable log-joint density function for
gradient-based inference.

The package addresses three limitations of existing tools:

\begin{enumerate}
\def\labelenumi{\arabic{enumi}.}
\tightlist
\item
  \textbf{greta} (Golding 2019) provides an elegant R-native DSL but
  requires Python, TensorFlow, and reticulate --- a fragile dependency
  chain.
\item
  \textbf{Stan/brms} (Carpenter et al. 2017; Bürkner 2017) requires
  learning a separate modelling language and managing a C++ toolchain.
\item
  \textbf{NIMBLE} (Valpine et al. 2017) uses BUGS-like syntax that
  differs from standard R.
\end{enumerate}

gretaR resolves these by providing pure R syntax with a torch backend
(no Python), optional Stan backend (no new language), and five inference
methods in a single package.

\subsection{Scope of this Document}\label{scope-of-this-document}

This document covers every major component of gretaR v0.1.0:

\begin{itemize}
\tightlist
\item
  The directed acyclic graph (DAG) and lazy evaluation system (Section
  2)
\item
  The 18 probability distributions and parameter transforms (Section 3)
\item
  Five inference engines: NUTS, HMC, ADVI, MAP, Laplace (Section 4)
\item
  Performance optimisations (Section 5)
\item
  The Stan code generation backend (Section 6)
\item
  The formula interface and random effects (Section 7)
\item
  Custom distributions and mixture models (Section 8)
\item
  Sparse matrix support (Section 9)
\item
  The unified output structure (Section 10)
\item
  Validation results (Section 11)
\end{itemize}

\section{Core Architecture: The Directed Acyclic
Graph}\label{core-architecture-the-directed-acyclic-graph}

\subsection{Lazy Evaluation and Node
Types}\label{lazy-evaluation-and-node-types}

gretaR uses lazy evaluation: operations on \texttt{gretaR\_array}
objects do not compute values immediately. Instead, they construct a DAG
of interconnected nodes. Each node is an R6 object
(\texttt{GretaRArray}) with fields:

\begin{itemize}
\tightlist
\item
  \textbf{\texttt{node\_type}}: one of \texttt{"data"},
  \texttt{"variable"}, or \texttt{"operation"}
\item
  \textbf{\texttt{value}}: a torch tensor (fixed for data; managed by
  the sampler for variables)
\item
  \textbf{\texttt{parents}}: character vector of parent node identifiers
\item
  \textbf{\texttt{operation}}: a closure that computes this node's value
  from its parents
\item
  \textbf{\texttt{distribution}}: a prior distribution (for variable
  nodes)
\item
  \textbf{\texttt{constraint}}: bounds on the parameter space
\item
  \textbf{\texttt{transform}}: a bijector mapping between constrained
  and unconstrained spaces
\end{itemize}

\subsection{Model Compilation}\label{model-compilation}

The \texttt{model()} function compiles the DAG into a differentiable
log-joint density function:

\begin{enumerate}
\def\labelenumi{\arabic{enumi}.}
\tightlist
\item
  \textbf{Identify free variables}: traverse the DAG to find all
  variable nodes
\item
  \textbf{Exclude likelihood templates}: variable nodes used on the
  right-hand side of \texttt{distribution(y)\ \textless{}-} are not free
  parameters
\item
  \textbf{Build parameter layout}: compute offsets for packing/unpacking
  a flat parameter vector
  \(\boldsymbol\theta_{\text{free}} \in \mathbb{R}^d\)
\item
  \textbf{Construct \texttt{log\_prob()}}: a function
  \(f: \mathbb{R}^d \to \mathbb{R}\) that computes:
\end{enumerate}

\[\log p(\boldsymbol\theta, \mathbf{y}) = \underbrace{\sum_{k=1}^{K} \log p(\theta_k \mid \text{prior}_k)}_{\text{priors}} + \underbrace{\sum_{k=1}^{K} \log \left| \frac{\partial T_k^{-1}}{\partial \theta_k^{\text{free}}} \right|}_{\text{Jacobian}} + \underbrace{\sum_{i=1}^{N} \log p(y_i \mid \boldsymbol\theta)}_{\text{likelihood}}\]

where \(T_k\) is the bijector for parameter \(k\).

\subsection{Gradient Computation}\label{gradient-computation}

Gradients are computed via torch's automatic differentiation:

\begin{Shaded}
\begin{Highlighting}[]
\NormalTok{grad\_log\_prob }\OtherTok{\textless{}{-}} \ControlFlowTok{function}\NormalTok{(model, theta\_free) \{}
\NormalTok{  theta\_free}\SpecialCharTok{$}\FunctionTok{requires\_grad\_}\NormalTok{(}\ConstantTok{TRUE}\NormalTok{)}
\NormalTok{  lp }\OtherTok{\textless{}{-}} \FunctionTok{log\_prob}\NormalTok{(model, theta\_free)}
\NormalTok{  grads }\OtherTok{\textless{}{-}} \FunctionTok{autograd\_grad}\NormalTok{(lp, theta\_free)}
  \FunctionTok{list}\NormalTok{(}\AttributeTok{lp =}\NormalTok{ lp}\SpecialCharTok{$}\FunctionTok{item}\NormalTok{(), }\AttributeTok{grad =}\NormalTok{ grads[[}\DecValTok{1}\NormalTok{]])}
\NormalTok{\}}
\end{Highlighting}
\end{Shaded}

The use of \texttt{autograd\_grad()} (rather than \texttt{backward()})
avoids gradient accumulation overhead and eliminates the need for
gradient zeroing between evaluations.

\section{Probability Distributions}\label{probability-distributions}

gretaR implements 18 probability distributions, each as an R6 class
inheriting from \texttt{GretaRDistribution}. Every distribution
provides:

\begin{itemize}
\tightlist
\item
  \texttt{log\_prob(x)}: differentiable log-density evaluated at \(x\)
  (returns a scalar torch tensor)
\item
  \texttt{sample(n)}: draw \(n\) samples (for prior predictive checks)
\item
  \texttt{constraint}: parameter space bounds (used to select the
  appropriate bijector)
\end{itemize}

\subsection{Continuous Distributions}\label{continuous-distributions}

\begin{longtable}[]{@{}
  >{\raggedright\arraybackslash}p{(\linewidth - 6\tabcolsep) * \real{0.2500}}
  >{\raggedright\arraybackslash}p{(\linewidth - 6\tabcolsep) * \real{0.4000}}
  >{\raggedright\arraybackslash}p{(\linewidth - 6\tabcolsep) * \real{0.1875}}
  >{\raggedright\arraybackslash}p{(\linewidth - 6\tabcolsep) * \real{0.1625}}@{}}
\caption{Continuous distributions. *Simplex, correlation, and PD
transforms are deferred to a future release.}\tabularnewline
\toprule\noalign{}
\begin{minipage}[b]{\linewidth}\raggedright
Distribution
\end{minipage} & \begin{minipage}[b]{\linewidth}\raggedright
Constructor
\end{minipage} & \begin{minipage}[b]{\linewidth}\raggedright
Support
\end{minipage} & \begin{minipage}[b]{\linewidth}\raggedright
Transform
\end{minipage} \\
\midrule\noalign{}
\endfirsthead
\toprule\noalign{}
\begin{minipage}[b]{\linewidth}\raggedright
Distribution
\end{minipage} & \begin{minipage}[b]{\linewidth}\raggedright
Constructor
\end{minipage} & \begin{minipage}[b]{\linewidth}\raggedright
Support
\end{minipage} & \begin{minipage}[b]{\linewidth}\raggedright
Transform
\end{minipage} \\
\midrule\noalign{}
\endhead
\bottomrule\noalign{}
\endlastfoot
Normal & \texttt{normal()} & \(\mathbb{R}\) & Identity \\
Half-Normal & \texttt{half\_normal()} & \(\mathbb{R}^+\) & Log \\
Half-Cauchy & \texttt{half\_cauchy()} & \(\mathbb{R}^+\) & Log \\
Student-\(t\) & \texttt{student\_t()} & \(\mathbb{R}\) & Identity \\
Cauchy & \texttt{cauchy()} & \(\mathbb{R}\) & Identity \\
Exponential & \texttt{exponential()} & \(\mathbb{R}^+\) & Log \\
Gamma & \texttt{gamma\_dist()} & \(\mathbb{R}^+\) & Log \\
Beta & \texttt{beta\_dist()} & \([0,1]\) & Logit \\
Log-Normal & \texttt{lognormal()} & \(\mathbb{R}^+\) & Log \\
Uniform & \texttt{uniform()} & \([a,b]\) & Scaled logit \\
Multivariate Normal & \texttt{multivariate\_normal()} & \(\mathbb{R}^k\)
& Identity \\
Dirichlet & \texttt{dirichlet()} & Simplex & Identity* \\
LKJ Correlation & \texttt{lkj\_correlation()} & Corr. matrix &
Identity* \\
Wishart & \texttt{wishart()} & PD matrix & Identity* \\
\end{longtable}

\subsection{Discrete Distributions}\label{discrete-distributions}

\begin{longtable}[]{@{}
  >{\raggedright\arraybackslash}p{(\linewidth - 4\tabcolsep) * \real{0.2571}}
  >{\raggedright\arraybackslash}p{(\linewidth - 4\tabcolsep) * \real{0.4286}}
  >{\raggedright\arraybackslash}p{(\linewidth - 4\tabcolsep) * \real{0.3143}}@{}}
\caption{Discrete distributions (used as likelihoods; not directly
sampled via HMC).}\tabularnewline
\toprule\noalign{}
\begin{minipage}[b]{\linewidth}\raggedright
Distribution
\end{minipage} & \begin{minipage}[b]{\linewidth}\raggedright
Constructor
\end{minipage} & \begin{minipage}[b]{\linewidth}\raggedright
Support
\end{minipage} \\
\midrule\noalign{}
\endfirsthead
\toprule\noalign{}
\begin{minipage}[b]{\linewidth}\raggedright
Distribution
\end{minipage} & \begin{minipage}[b]{\linewidth}\raggedright
Constructor
\end{minipage} & \begin{minipage}[b]{\linewidth}\raggedright
Support
\end{minipage} \\
\midrule\noalign{}
\endhead
\bottomrule\noalign{}
\endlastfoot
Bernoulli & \texttt{bernoulli()} & \(\{0,1\}\) \\
Binomial & \texttt{binomial()} & \(\{0,\ldots,n\}\) \\
Poisson & \texttt{poisson\_dist()} & \(\mathbb{Z}_{\geq 0}\) \\
Negative Binomial & \texttt{negative\_binomial()} &
\(\mathbb{Z}_{\geq 0}\) \\
\end{longtable}

\subsection{Parameter Transforms
(Bijectors)}\label{parameter-transforms-bijectors}

Six bijectors map between constrained parameter spaces and the
unconstrained real line, as required by gradient-based samplers:

{\def\LTcaptype{none} % do not increment counter
\begin{longtable}[]{@{}
  >{\raggedright\arraybackslash}p{(\linewidth - 8\tabcolsep) * \real{0.2000}}
  >{\raggedright\arraybackslash}p{(\linewidth - 8\tabcolsep) * \real{0.2000}}
  >{\raggedright\arraybackslash}p{(\linewidth - 8\tabcolsep) * \real{0.2000}}
  >{\raggedright\arraybackslash}p{(\linewidth - 8\tabcolsep) * \real{0.2000}}
  >{\raggedright\arraybackslash}p{(\linewidth - 8\tabcolsep) * \real{0.2000}}@{}}
\toprule\noalign{}
\begin{minipage}[b]{\linewidth}\raggedright
Bijector
\end{minipage} & \begin{minipage}[b]{\linewidth}\raggedright
Domain
\end{minipage} & \begin{minipage}[b]{\linewidth}\raggedright
Forward \(T(x)\)
\end{minipage} & \begin{minipage}[b]{\linewidth}\raggedright
Inverse \(T^{-1}(y)\)
\end{minipage} & \begin{minipage}[b]{\linewidth}\raggedright
\(\log|J|\)
\end{minipage} \\
\midrule\noalign{}
\endhead
\bottomrule\noalign{}
\endlastfoot
Identity & \(\mathbb{R}\) & \(y = x\) & \(x = y\) & \(0\) \\
Log & \(\mathbb{R}^+\) & \(y = \log x\) & \(x = e^y\) & \(\sum y_i\) \\
Logit & \((0,1)\) & \(y = \log\frac{x}{1-x}\) & \(x = \sigma(y)\) &
\(\sum[\log\sigma(y) + \log(1-\sigma(y))]\) \\
Scaled logit & \((a,b)\) & Shifted logit & Shifted sigmoid & Includes
range adjustment \\
Softplus & \(\mathbb{R}^+\) & \(y = \log(e^x - 1)\) &
\(x = \text{softplus}(y)\) & \(\sum\log\sigma(y)\) \\
Lower bound & \((L, \infty)\) & \(y = \log(x - L)\) & \(x = e^y + L\) &
\(\sum y_i\) \\
\end{longtable}
}

The \texttt{select\_transform()} function automatically chooses the
bijector based on the distribution's constraint specification.

\section{Inference Engines}\label{inference-engines}

\subsection{No-U-Turn Sampler (NUTS)}\label{no-u-turn-sampler-nuts}

The primary inference method implements the NUTS algorithm (Hoffman and
Gelman 2014) with the following components:

\textbf{Leapfrog integrator.} Given position \(\boldsymbol\theta\) and
momentum \(\mathbf{p}\), one step with step size \(\epsilon\) and mass
matrix \(\mathbf{M}\):

\begin{align}
\mathbf{p} &\leftarrow \mathbf{p} + \frac{\epsilon}{2} \nabla_\theta \log p(\boldsymbol\theta, \mathbf{y}) \\
\boldsymbol\theta &\leftarrow \boldsymbol\theta + \epsilon \mathbf{M}^{-1} \mathbf{p} \\
\mathbf{p} &\leftarrow \mathbf{p} + \frac{\epsilon}{2} \nabla_\theta \log p(\boldsymbol\theta, \mathbf{y})
\end{align}

\textbf{Tree doubling.} The trajectory is extended by doubling in a
randomly chosen direction. At each depth \(j\), \(2^j\) leapfrog steps
are evaluated. Growth stops when a U-turn is detected:
\(\Delta\boldsymbol\theta \cdot \mathbf{p}_{\text{left}} < 0\) or
\(\Delta\boldsymbol\theta \cdot \mathbf{p}_{\text{right}} < 0\).

\textbf{Windowed warmup.} Adaptation follows a three-phase strategy
inspired by Stan (Carpenter et al. 2017):

\begin{itemize}
\tightlist
\item
  Phase 1 (0--15\%): step-size adaptation via dual averaging (Nesterov
  2009)
\item
  Phase 2 (15--90\%): collect samples for diagonal mass matrix
  estimation
\item
  Phase 3 (90--100\%): re-adapt step size with the learned mass matrix
\end{itemize}

This phased approach prevents the mass matrix update from invalidating
the previously adapted step size.

\textbf{Implementation note.} The leapfrog operates on plain R numeric
vectors (not torch tensors) to avoid scoping issues with torch's
\texttt{with\_no\_grad()} in R. Only the gradient evaluation touches
torch.

\subsection{Hamiltonian Monte Carlo
(HMC)}\label{hamiltonian-monte-carlo-hmc}

Static HMC with a fixed number of leapfrog steps per iteration. Uses
identical warmup adaptation to NUTS. Recommended for debugging or when
NUTS tree depth is a concern.

\subsection{Automatic Differentiation Variational Inference
(ADVI)}\label{automatic-differentiation-variational-inference-advi}

Implements the ADVI framework of Kucukelbir et al. (2017) with two
variational families:

\textbf{Mean-field:}
\(q(\boldsymbol\theta) = \mathcal{N}(\boldsymbol\mu, \text{diag}(\boldsymbol\sigma^2))\)

\textbf{Full-rank:}
\(q(\boldsymbol\theta) = \mathcal{N}(\boldsymbol\mu, \mathbf{L}\mathbf{L}^\top)\)

The ELBO is optimised via the reparameterisation trick:

\[\text{ELBO} = \mathbb{E}_{q}[\log p(\mathbf{y}, \boldsymbol\theta)] + \mathcal{H}[q]\]

where
\(\boldsymbol\theta = \boldsymbol\mu + \boldsymbol\sigma \odot \boldsymbol\epsilon\)
and \(\boldsymbol\epsilon \sim \mathcal{N}(\mathbf{0}, \mathbf{I})\).

Convergence is assessed by monitoring the relative change in the median
ELBO over a rolling window.

\subsection{MAP Estimation}\label{map-estimation}

Maximum a posteriori estimation via the Adam optimiser (Kingma and Ba
2015). The initial point is found by gradient ascent toward the
posterior mode (200 steps with learning rate 0.1), providing robust
initialisation even for multimodal posteriors.

\subsection{Laplace Approximation}\label{laplace-approximation}

Approximates the posterior as a multivariate normal centred at the MAP:

\[p(\boldsymbol\theta \mid \mathbf{y}) \approx \mathcal{N}\left(\hat{\boldsymbol\theta}_{\text{MAP}}, \left[-\nabla^2 \log p(\hat{\boldsymbol\theta}_{\text{MAP}}, \mathbf{y})\right]^{-1}\right)\]

The Hessian is computed via symmetric finite differences on the
gradient. Eigenvalue clamping ensures positive definiteness. The log
marginal likelihood is approximated as:

\[\log p(\mathbf{y}) \approx \log p(\hat{\boldsymbol\theta}, \mathbf{y}) + \frac{d}{2}\log(2\pi) + \frac{1}{2}\log|\boldsymbol\Sigma|\]

\section{Performance Optimisations}\label{performance-optimisations}

Three optimisations reduce the torch backend's per-gradient-evaluation
cost:

\subsection{Compiled Log-Probability
Function}\label{compiled-log-probability-function}

At model compilation time, \texttt{compile\_log\_prob()} pre-extracts
all parameter layout information (offsets, dimensions, transforms,
distribution references) into a closure. This eliminates R6 method
dispatch and list lookups during the hot path of gradient evaluation.

\subsection{\texorpdfstring{\texttt{autograd\_grad()} vs
\texttt{backward()}}{autograd\_grad() vs backward()}}\label{autograd_grad-vs-backward}

The standard \texttt{backward()} accumulates gradients into the
\texttt{\$grad} attribute, requiring subsequent \texttt{\$clone()} and
\texttt{\$detach()\$cpu()} calls. Using \texttt{autograd\_grad()}
returns gradients directly, reducing per-call overhead by approximately
27\%.

\subsection{JIT Tracing}\label{jit-tracing}

The compiled log-probability function is optionally traced via
\texttt{torch::jit\_trace()}, which fuses torch operations into an
optimised graph. Tracing succeeds for models without data-dependent
control flow; a safe fallback to the untraced compiled function is used
otherwise.

\textbf{Combined effect:} 3--4\(\times\) speedup over the initial
implementation, reducing the torch--Stan gap from 50--64\(\times\) to
7--16\(\times\) for small models.

\subsubsection{Approaches Evaluated and
Rejected}\label{approaches-evaluated-and-rejected}

{\def\LTcaptype{none} % do not increment counter
\begin{longtable}[]{@{}
  >{\raggedright\arraybackslash}p{(\linewidth - 4\tabcolsep) * \real{0.3333}}
  >{\raggedright\arraybackslash}p{(\linewidth - 4\tabcolsep) * \real{0.3333}}
  >{\raggedright\arraybackslash}p{(\linewidth - 4\tabcolsep) * \real{0.3333}}@{}}
\toprule\noalign{}
\begin{minipage}[b]{\linewidth}\raggedright
Approach
\end{minipage} & \begin{minipage}[b]{\linewidth}\raggedright
Finding
\end{minipage} & \begin{minipage}[b]{\linewidth}\raggedright
Decision
\end{minipage} \\
\midrule\noalign{}
\endhead
\bottomrule\noalign{}
\endlastfoot
GPU/MPS acceleration & 5.6\(\times\) \emph{slower} for small tensors due
to transfer latency & Rejected \\
C++ leapfrog loop & R-level loop overhead is \$\textless\$1\% of total
time & Rejected \\
In-place tensor operations & Higher dispatch cost than allocation for
small tensors & Rejected \\
Batched chains (vectorised HMC) & NUTS trees vary per chain, wasting
compute & Deferred \\
\end{longtable}
}

\section{Stan Backend}\label{stan-backend}

\subsection{Motivation}\label{motivation}

The torch backend executes the DAG interpretively: each gradient
evaluation traverses R-level data structures and dispatches torch
operations individually. Stan compiles the entire model to optimised C++
with static automatic differentiation, yielding 32--152\(\times\) faster
sampling on hierarchical models.

gretaR's Stan backend provides this speed advantage while preserving the
R-native DSL:

\begin{Shaded}
\begin{Highlighting}[]
\CommentTok{\# Same model definition}
\NormalTok{mu }\OtherTok{\textless{}{-}} \FunctionTok{normal}\NormalTok{(}\DecValTok{0}\NormalTok{, }\DecValTok{10}\NormalTok{)}
\NormalTok{sigma }\OtherTok{\textless{}{-}} \FunctionTok{half\_cauchy}\NormalTok{(}\DecValTok{2}\NormalTok{)}
\FunctionTok{distribution}\NormalTok{(y) }\OtherTok{\textless{}{-}} \FunctionTok{normal}\NormalTok{(mu, sigma)}
\NormalTok{m }\OtherTok{\textless{}{-}} \FunctionTok{model}\NormalTok{(mu, sigma)}

\CommentTok{\# Choose backend at inference time}
\NormalTok{fit\_torch }\OtherTok{\textless{}{-}} \FunctionTok{mcmc}\NormalTok{(m, }\AttributeTok{backend =} \StringTok{"torch"}\NormalTok{)  }\CommentTok{\# native torch NUTS}
\NormalTok{fit\_stan  }\OtherTok{\textless{}{-}} \FunctionTok{mcmc}\NormalTok{(m, }\AttributeTok{backend =} \StringTok{"stan"}\NormalTok{)   }\CommentTok{\# compiled Stan NUTS}
\end{Highlighting}
\end{Shaded}

\subsection{Code Generation}\label{code-generation}

\texttt{compile\_to\_stan()} walks the compiled model's DAG and emits
valid Stan code:

\begin{enumerate}
\def\labelenumi{\arabic{enumi}.}
\tightlist
\item
  \textbf{Data block}: declares observed data with dimensions inferred
  from tensor shapes
\item
  \textbf{Parameters block}: declares free parameters with constraints
  (e.g., \texttt{real\textless{}lower=0\textgreater{}})
\item
  \textbf{Model block}: emits prior statements and likelihood statements
\end{enumerate}

Each gretaR distribution maps to its Stan equivalent (e.g.,
\texttt{half\_cauchy(scale)} \(\to\) \texttt{cauchy(0,\ scale)} with
\texttt{\textless{}lower=0\textgreater{}} bound).

\subsection{Benchmark Results}\label{benchmark-results}

\begin{longtable}[]{@{}lrrl@{}}
\caption{Wall-clock time (500 warmup + 500 samples, 2 chains). Parameter
recovery matches to 3 significant figures across
backends.}\tabularnewline
\toprule\noalign{}
Model & Torch (s) & Stan (s) & Speedup \\
\midrule\noalign{}
\endfirsthead
\toprule\noalign{}
Model & Torch (s) & Stan (s) & Speedup \\
\midrule\noalign{}
\endhead
\bottomrule\noalign{}
\endlastfoot
Intercept-only (2p) & 351.1 & 3.7 & 95x \\
Linear regression (3p) & 2.9 & 3.4 & 1x \\
Random intercept LMM (23p) & 466.5 & 4.6 & 101x \\
Crossed random effects (42p) & 1154.5 & 7.6 & 152x \\
Logistic GLMM (17p) & 423.1 & 8.6 & 49x \\
Poisson GLMM (22p) & 441.8 & 5.9 & 75x \\
Random slopes (35p) & 208.8 & 6.5 & 32x \\
\end{longtable}

\section{Formula Interface}\label{formula-interface}

\subsection{Basic Usage}\label{basic-usage}

\texttt{gretaR\_glm()} provides a high-level interface using standard R
formula syntax:

\begin{Shaded}
\begin{Highlighting}[]
\NormalTok{fit }\OtherTok{\textless{}{-}} \FunctionTok{gretaR\_glm}\NormalTok{(}
\NormalTok{  Sepal.Length }\SpecialCharTok{\textasciitilde{}}\NormalTok{ Sepal.Width }\SpecialCharTok{+}\NormalTok{ Petal.Length,}
  \AttributeTok{data =}\NormalTok{ iris,}
  \AttributeTok{family =} \StringTok{"gaussian"}
\NormalTok{)}
\end{Highlighting}
\end{Shaded}

Internally, \texttt{model.matrix()} constructs the design matrix;
default priors are assigned (\(\beta \sim \mathcal{N}(0, 5)\),
\(\sigma \sim \text{Half-Cauchy}(2)\)); and the model is compiled and
sampled via the standard gretaR pipeline.

\subsection{Random Effects}\label{random-effects}

lme4-style random effects are parsed via regex (no lme4 dependency
required):

\begin{Shaded}
\begin{Highlighting}[]
\NormalTok{fit }\OtherTok{\textless{}{-}} \FunctionTok{gretaR\_glm}\NormalTok{(}
\NormalTok{  y }\SpecialCharTok{\textasciitilde{}}\NormalTok{ x }\SpecialCharTok{+}\NormalTok{ (}\DecValTok{1} \SpecialCharTok{|}\NormalTok{ group),        }\CommentTok{\# random intercepts}
  \AttributeTok{data =}\NormalTok{ dat,}
  \AttributeTok{family =} \StringTok{"gaussian"}
\NormalTok{)}
\end{Highlighting}
\end{Shaded}

Supported patterns:

\begin{itemize}
\tightlist
\item
  \texttt{(1\ \textbar{}\ group)} --- random intercepts
\item
  \texttt{(x\ \textbar{}\ group)} --- correlated random intercepts and
  slopes
\item
  \texttt{(0\ +\ x\ \textbar{}\ group)} --- random slopes only
\item
  Multiple terms:
  \texttt{(1\ \textbar{}\ site)\ +\ (1\ \textbar{}\ year)}
\end{itemize}

All random effects use non-centred parameterisation by default:
\(\alpha_j = \tau \cdot z_j\) where \(z_j \sim \mathcal{N}(0, 1)\) and
\(\tau \sim \text{Half-Cauchy}(2)\). This avoids the funnel geometry
that degrades HMC performance in centred parameterisations (Betancourt
and Girolami 2015).

\subsection{Formula Style Detection}\label{formula-style-detection}

The parser auto-detects formula styles:

{\def\LTcaptype{none} % do not increment counter
\begin{longtable}[]{@{}lll@{}}
\toprule\noalign{}
Pattern & Detected as & Action \\
\midrule\noalign{}
\endhead
\bottomrule\noalign{}
\endlastfoot
\texttt{y\ \textasciitilde{}\ x1\ +\ x2} & base R & Standard GLM via
\texttt{model.matrix()} \\
\texttt{y\ \textasciitilde{}\ x\ +\ (1\textbackslash{}\textbar{}group)}
& lme4 & Mixed model with non-centred RE \\
\texttt{y\ \textasciitilde{}\ s(x)} & mgcv & Error (not yet
supported) \\
\end{longtable}
}

\section{Custom Distributions and Mixture
Models}\label{custom-distributions-and-mixture-models}

\subsection{Custom Distributions}\label{custom-distributions}

Users can define distributions with arbitrary torch-differentiable
log-probability functions:

\begin{Shaded}
\begin{Highlighting}[]
\NormalTok{x }\OtherTok{\textless{}{-}} \FunctionTok{custom\_distribution}\NormalTok{(}
  \AttributeTok{log\_prob\_fn =} \ControlFlowTok{function}\NormalTok{(x) }\SpecialCharTok{{-}}\FunctionTok{torch\_sum}\NormalTok{(}\FunctionTok{torch\_abs}\NormalTok{(x)),}
  \AttributeTok{constraint =} \FunctionTok{list}\NormalTok{(}\AttributeTok{lower =} \SpecialCharTok{{-}}\ConstantTok{Inf}\NormalTok{, }\AttributeTok{upper =} \ConstantTok{Inf}\NormalTok{),}
  \AttributeTok{name =} \StringTok{"laplace"}
\NormalTok{)}
\end{Highlighting}
\end{Shaded}

The supplied function must accept a torch tensor and return a scalar
torch tensor. It must be differentiable via torch autograd.

\subsection{Mixture Models}\label{mixture-models}

Finite mixture models are supported via the log-sum-exp trick:

\[\log p(x \mid \boldsymbol\pi, \boldsymbol\Theta) = \log \sum_{k=1}^{K} \pi_k \, p_k(x \mid \theta_k) = \text{logsumexp}_k \left[\log \pi_k + \log p_k(x \mid \theta_k)\right]\]

This marginalises over the discrete component indicator, enabling
gradient-based inference:

\begin{Shaded}
\begin{Highlighting}[]
\NormalTok{w }\OtherTok{\textless{}{-}} \FunctionTok{dirichlet}\NormalTok{(}\FunctionTok{c}\NormalTok{(}\DecValTok{1}\NormalTok{, }\DecValTok{1}\NormalTok{))}
\NormalTok{mu1 }\OtherTok{\textless{}{-}} \FunctionTok{normal}\NormalTok{(}\SpecialCharTok{{-}}\DecValTok{2}\NormalTok{, }\DecValTok{1}\NormalTok{); mu2 }\OtherTok{\textless{}{-}} \FunctionTok{normal}\NormalTok{(}\DecValTok{2}\NormalTok{, }\DecValTok{1}\NormalTok{)}
\NormalTok{sigma }\OtherTok{\textless{}{-}} \FunctionTok{half\_cauchy}\NormalTok{(}\DecValTok{1}\NormalTok{)}

\NormalTok{mix }\OtherTok{\textless{}{-}} \FunctionTok{mixture}\NormalTok{(}
  \AttributeTok{distributions =} \FunctionTok{list}\NormalTok{(}\FunctionTok{normal}\NormalTok{(mu1, sigma), }\FunctionTok{normal}\NormalTok{(mu2, sigma)),}
  \AttributeTok{weights =}\NormalTok{ w}
\NormalTok{)}
\FunctionTok{distribution}\NormalTok{(y) }\OtherTok{\textless{}{-}}\NormalTok{ mix}
\end{Highlighting}
\end{Shaded}

\section{Sparse Matrix Support}\label{sparse-matrix-support}

Large, sparse design matrices (common in genomics, NLP, and spatial
modelling) are handled via the Matrix package (Bates et al. 2024). The
\texttt{as\_data()} function dispatches on \texttt{sparseMatrix}
objects:

\begin{Shaded}
\begin{Highlighting}[]
\FunctionTok{library}\NormalTok{(Matrix)}
\NormalTok{X\_sparse }\OtherTok{\textless{}{-}} \FunctionTok{sparseMatrix}\NormalTok{(}\AttributeTok{i =}\NormalTok{ ..., }\AttributeTok{j =}\NormalTok{ ..., }\AttributeTok{x =}\NormalTok{ ...)}
\NormalTok{X }\OtherTok{\textless{}{-}} \FunctionTok{as\_data}\NormalTok{(X\_sparse)  }\CommentTok{\# stored as torch sparse COO tensor}
\NormalTok{mu }\OtherTok{\textless{}{-}}\NormalTok{ X }\SpecialCharTok{\%*\%}\NormalTok{ beta         }\CommentTok{\# sparse{-}aware matrix multiplication}
\end{Highlighting}
\end{Shaded}

The Matrix package is a recommended R package shipped with every R
installation, imposing zero additional installation burden. Sparse
\(\times\) dense matrix multiplication uses \texttt{torch\_mm()}, which
handles the sparse layout natively.

\section{Unified Output Structure}\label{unified-output-structure}

All inference functions return a \texttt{gretaR\_fit} S3 object with
consistent structure:

\begin{Shaded}
\begin{Highlighting}[]
\NormalTok{fit}\SpecialCharTok{$}\NormalTok{draws        }\CommentTok{\# posterior::draws\_array (NULL for MAP/Laplace)}
\NormalTok{fit}\SpecialCharTok{$}\NormalTok{model        }\CommentTok{\# compiled gretaR\_model}
\NormalTok{fit}\SpecialCharTok{$}\NormalTok{summary      }\CommentTok{\# tibble: mean, sd, quantiles, R{-}hat, ESS}
\NormalTok{fit}\SpecialCharTok{$}\NormalTok{convergence  }\CommentTok{\# list: n\_eff, rhat, max\_rhat, min\_ess, n\_divergences}
\NormalTok{fit}\SpecialCharTok{$}\NormalTok{call\_info    }\CommentTok{\# original arguments for reproducibility}
\NormalTok{fit}\SpecialCharTok{$}\NormalTok{run\_time     }\CommentTok{\# elapsed seconds}
\NormalTok{fit}\SpecialCharTok{$}\NormalTok{method       }\CommentTok{\# "nuts", "hmc", "vi", "map", "laplace", "stan"}
\NormalTok{fit}\SpecialCharTok{$}\NormalTok{par          }\CommentTok{\# point estimates (all methods)}
\end{Highlighting}
\end{Shaded}

S3 methods provide a consistent interface:

\begin{itemize}
\tightlist
\item
  \texttt{print(fit)} --- concise summary with convergence diagnostics
\item
  \texttt{summary(fit)} --- full posterior table via
  \texttt{posterior::summarise\_draws()}
\item
  \texttt{coef(fit)} --- named vector of posterior means or MAP
  estimates
\item
  \texttt{plot(fit,\ type)} --- diagnostic plots via bayesplot
  (\texttt{"trace"}, \texttt{"density"}, \texttt{"pairs"},
  \texttt{"rhat"}, \texttt{"neff"})
\end{itemize}

\section{Validation}\label{validation}

gretaR was validated against CmdStan (Carpenter et al. 2017) on 10
benchmark models. Parameter estimates agree to 2--3 significant figures
across all benchmarks. The validation suite is included in
\texttt{inst/validation/} and covers:

\begin{itemize}
\tightlist
\item
  Normal mean estimation (known and unknown variance)
\item
  Linear and multiple regression
\item
  Logistic and Poisson regression
\item
  Gamma GLM and robust regression (Student-\(t\) errors)
\item
  Beta regression
\item
  Hierarchical random intercepts (non-centred parameterisation)
\end{itemize}

\subsection{Correctness Criteria}\label{correctness-criteria}

A model is considered validated when:

\begin{itemize}
\tightlist
\item
  \(\hat{R} < 1.05\) for all parameters
\item
  Bulk ESS \(> 400\) for all parameters
\item
  Posterior means within MCMC uncertainty of Stan reference values
\item
  Zero post-warmup divergent transitions
\end{itemize}

\section{Potential for Further
Development}\label{potential-for-further-development}

\subsection{Near-Term (Phase 3)}\label{near-term-phase-3}

\begin{itemize}
\tightlist
\item
  \textbf{mgcv smooth terms} (\texttt{s()}, \texttt{te()}) via basis
  function expansion
\item
  \textbf{Simplex and correlation transforms} for Dirichlet and LKJ
  sampling
\item
  \textbf{Wishart/inverse-Wishart sampling} via Bartlett decomposition
\item
  \textbf{Improved Stan code generator} with operation-type metadata on
  DAG nodes
\end{itemize}

\subsection{Medium-Term}\label{medium-term}

\begin{itemize}
\tightlist
\item
  \textbf{Gaussian Process module} (\texttt{gretaR.gp}) as an extension
  package
\item
  \textbf{ODE-based models} using torch ODE solvers
\item
  \textbf{Expectation propagation} for approximate inference
\item
  \textbf{S7 class system} migration (when mature)
\end{itemize}

\subsection{Long-Term}\label{long-term}

\begin{itemize}
\tightlist
\item
  \textbf{GPU-accelerated inference} for large latent variable models
\item
  \textbf{Batched chain evaluation} for static HMC on GPU
\item
  \textbf{Community extension framework} with templates and
  documentation
\item
  \textbf{CRAN submission} and JOSS publication
\end{itemize}

\section*{References}\label{references}
\addcontentsline{toc}{section}{References}

\protect\phantomsection\label{refs}
\begin{CSLReferences}{1}{1}
\bibitem[\citeproctext]{ref-matrix}
Bates, Douglas, Martin Maechler, and Mikael Jagan. 2024. \emph{Matrix:
Sparse and Dense Matrix Classes and Methods}.
\url{https://CRAN.R-project.org/package=Matrix}.

\bibitem[\citeproctext]{ref-betancourt2015}
Betancourt, Michael, and Mark Girolami. 2015. {``Hamiltonian {M}onte
{C}arlo for Hierarchical Models.''} \emph{Current Trends in Bayesian
Methodology with Applications}, 79--101.

\bibitem[\citeproctext]{ref-brms}
Bürkner, Paul-Christian. 2017. {``{brms}: An {R} Package for {B}ayesian
Multilevel Models Using {Stan}.''} \emph{Journal of Statistical
Software} 80 (1): 1--28. \url{https://doi.org/10.18637/jss.v080.i01}.

\bibitem[\citeproctext]{ref-stan}
Carpenter, Bob, Andrew Gelman, Matthew D. Hoffman, et al. 2017. {``Stan:
A Probabilistic Programming Language.''} \emph{Journal of Statistical
Software} 76 (1): 1--32. \url{https://doi.org/10.18637/jss.v076.i01}.

\bibitem[\citeproctext]{ref-greta}
Golding, Nick. 2019. {``Greta: Simple and Scalable Statistical Modelling
in {R}.''} \emph{Journal of Open Source Software} 4 (40): 1601.
\url{https://doi.org/10.21105/joss.01601}.

\bibitem[\citeproctext]{ref-nuts}
Hoffman, Matthew D., and Andrew Gelman. 2014. {``The No-{U}-Turn
Sampler: Adaptively Setting Path Lengths in {H}amiltonian {M}onte
{C}arlo.''} \emph{Journal of Machine Learning Research} 15: 1593--623.

\bibitem[\citeproctext]{ref-kingma2015adam}
Kingma, Diederik P., and Jimmy Ba. 2015. {``Adam: A Method for
Stochastic Optimization.''} \emph{arXiv Preprint arXiv:1412.6980}.

\bibitem[\citeproctext]{ref-advi}
Kucukelbir, Alp, Dustin Tran, Rajesh Ranganath, Andrew Gelman, and David
M. Blei. 2017. {``Automatic Differentiation Variational Inference.''}
\emph{Journal of Machine Learning Research} 18 (14): 1--45.

\bibitem[\citeproctext]{ref-nesterov2009}
Nesterov, Yurii. 2009. {``Primal-Dual Subgradient Methods for Convex
Problems.''} \emph{Mathematical Programming} 120 (1): 221--59.
\url{https://doi.org/10.1007/s10107-007-0149-x}.

\bibitem[\citeproctext]{ref-nimble}
Valpine, Perry de, Daniel Turek, Christopher J. Paciorek, Clifford
Anderson-Bergman, Duncan Temple Lang, and Rastislav Bodik. 2017.
{``Programming with Models: Writing Statistical Algorithms for General
Model Structures with {NIMBLE}.''} \emph{Journal of Computational and
Graphical Statistics} 26 (2): 403--13.
\url{https://doi.org/10.1080/10618600.2016.1172487}.

\end{CSLReferences}

\end{document}
